\documentclass[12pt]{article}
\usepackage[utf8]{inputenc}
\usepackage[T1]{fontenc}
\usepackage{hyperref}
\title{Fasu Language Resources for K-12 Curriculum Development}
\author{Compiled by Claude Sonnet 5.5 (claude-sonnet-5-5) for Andr\'as Kornai}
\date{}
\begin{document}
\maketitle

\section{Introduction}
Fasu (ISO 639-3 \texttt{faa}), also called Namo Me (``real speech''), is a Kutubuan language spoken southwest of Lake Kutubu, Southern Highlands, Papua New Guinea. The figure of about 1,200 speakers dates from 1981; no current figure was found. ``Fasu'' properly names one clan of Namo Me speakers. Its inclusion in Trans-New Guinea is disputed.

\section{Primary Educational Resources}
No primers, textbooks or school programmes were found in the search results.

\section{Linguistic and Reference Materials}
\begin{itemize}
\item Fasu language (Wikipedia); it cites Loeweke and May's 1980 grammar and 1981 Fasu--English dictionary (SIL), which were not themselves located online in this search: \url{https://en.wikipedia.org/wiki/Fasu_language}
\item Kutubuan languages (Wikipedia): \url{https://en.wikipedia.org/wiki/Kutubuan_languages}
\item Lake Kutubu languages overview (New Guinea World, University of Zurich): \url{https://newguineaworld.linguistik.uzh.ch/families/papuan-gulf/kikori-river/lake-kutubu/start}
\item Namo Me (New Guinea World, University of Zurich): \url{https://newguineaworld.linguistik.uzh.ch/families/papuan-gulf/kikori-river/lake-kutubu/namo-me}
\end{itemize}

\section{Religious and Community Resources}
None found.

\section{Digital and Audio Resources}
None found.

\section{Curriculum Development Context}
Documentation is mainly older SIL linguistic work; no school materials were found. The neighbouring Foi language is a different language and should not be conflated.

\end{document}